\pagestyle{empty}
\setlength{\parindent}{0cm}

{\centering
    \textbf{Національний технічний університет України}

    \textbf{\quotes{Київський політехнічний інститут імені Ігоря Сікорського}}

    \textbf{Факультет прикладної математики}

    \textbf{Кафедра прикладної математики}

}

{\raggedright
Рівень вищої освіти --- другий (магістерський)

Спеціальність --- 113~Прикладна математика

Освітньо-наукова програма \quotes{Наука про дані та математичне моделювання}

}

\addvspace{12pt}

\begin{flushright}
	\renewcommand{\arraystretch}{0.8}
	\begin{tabular}{l}
		\MakeUppercase{Затверджую} \\
		Завідувач кафедри \\
		\_\_\_\_\_\_\_\_\_\_\_\_\_~Данило ТАВРОВ \\
		\quotes{\_\_\_\_}~\_\_\_\_\_\_\_\_\_\_\_\_\_~\the\year~р. \\
	\end{tabular}
\end{flushright}

{\centering
	\textbf{\MakeUppercase{Завдання}}
	
	\textbf{на магістерську дисертацію студенту}

	\reportAuthorGen

}

1. Тема дисертації: \quotes{\reportTitle},
керівник роботи \supervisorFio, \supervisorRegalia, затверджені наказом по університету \reportOrder.

2. Термін подання студентом дисертації: \applicationDate

3. Вихідні дані до роботи: вихідними даними є музичні аудіозаписи у звуковому форматі з відкритого датасету Isophonics, що містить акордові анотації для творів The Beatles у форматі часових міток, а також стандартизовані метрики оцінювання MIREX.

4. Зміст роботи: аналіз існуючих підходів до автоматичного гармонічного аналізу музики; дослідження методів цифрової обробки аудіосигналів (STFT, CQT) та виділення хроматичних ознак; розробка математичної моделі послідовнісної класифікації гармонічних функцій акордів; реалізація та навчання моделі на основі двонапрямленої рекурентної нейронної мережі (BiLSTM); експериментальне оцінювання якості запропонованого методу та порівняння з базовими підходами (HMM, SVM).

5. Перелік ілюстративного матеріалу: концептуальна схема взаємодії компонентів системи автоматичного гармонічного аналізу; графіки динаміки навчання моделі; таблиця порівняння методів автоматичного гармонічного аналізу; таблиця результатів оцінювання моделі на тестовій вибірці.

6. Дата видачі завдання: \assignmentDate

\begin{center}

Календарний план

\addvspace{12pt}

\DefTblrTemplate{firsthead, middlehead, lasthead, firstfoot, middlefoot, lastfoot}{default}{}

\begin{longtblr}{
    hlines,
    vlines,
    colspec = {Q[0.1\textwidth, c] Q[0.4\textwidth, l] Q[0.25\textwidth, c] Q[0.15\textwidth, c]},
    cells = {cmd={\linespread{1.2}\selectfont}},
    rowhead = {1},
    rows = {valign=m},
    row{1} = {halign=c},
}
\No\ з/п & Назва етапів виконання магістерської дисертації & Термін виконання етапів роботи & Примітка \\

1 & Узгодження теми дисертації з науковим керівником, огляд наукових публікацій за тематикою дослідження & ХХ 2025 р. & \\
2 & Аналіз існуючих підходів до автоматичного гармонічного аналізу музики, написання розділу огляду & ХХ 2025 р. & \\
3 & Дослідження методів цифрової обробки аудіосигналів, розробка математичної моделі & ХХ 2025 р. & \\
4 & Реалізація модулів попередньої обробки аудіосигналу та виділення хроматичних ознак & ХХ 2026 р. & \\
5 & Розробка та навчання моделі BiLSTM для послідовнісної класифікації гармонічних функцій & ХХ 2026 р. & \\
6 & Проведення експериментального дослідження, порівняння з базовими методами (HMM, SVM) & ХХ 2026 р. & \\
7 & Написання розділів дисертації, підготовка висновків та списку використаних джерел & ХХ 2026 р. & \\
8 & Остаточне оформлення дисертації, підготовка презентації & ХХ 2026 р. & \\
9 & Подання повної версії магістерської дисертації & ХХ 2026 р. & \\

\end{longtblr}

\vspace{-0.5cm}

\begin{tblr}{
    colspec = {Q[7cm, l] X[c] X[r]},
    cells = {cmd={\linespread{1.2}\selectfont}},
    column{1} = {leftsep=0pt},
    column{Z} = {rightsep=0pt},
}

Студент &  \_\_\_\_\_\_\_\_\_\_\_\_\_ & \reportAuthorShort \\
Керівник роботи & \_\_\_\_\_\_\_\_\_\_\_\_\_ & \supervisorFioShort \\

\end{tblr}

\end{center}

\clearpage
\setlength{\parindent}{1.25cm}